\section{Related Work}
\paragraph{Sequential Recommendation.}
Sequential recommendation models user history to predict future behaviors. GRU4Rec~\cite{hidasi2015session} pioneered RNNs for temporal modeling. SASRec~\cite{kang2018self} adapts self-attention with causal masking, whereas BERT4Rec~\cite{sun2019bert4rec} employs bidirectional objectives to deepen context understanding. Recent advances optimize this backbone. LightSANs~\cite{fan2021lighter} introduces low-rank decomposed attention for linear scalability, and FEARec~\cite{du2023frequency} leverages frequency domain learning for multi-scale information. However, these methods focus on fitting historical preferences and fail to shift user preferences.

\paragraph{Proactive Recommendation.}
Proactive recommendations aim to shift user preferences toward target items. IRN~\cite{zhu2023influential} introduces a Transformer-based supervised method with a Personalized Impressionability Mask to model user receptiveness. IPG~\cite{bi2024proactive} selects intermediate items by jointly evaluating local feasibility and guidance effectiveness via predefined heuristics, while ITMPRec~\cite{lian2025itmprec} further incorporates intention-level features for finer-grained characterization. More recently, LLM-IPP~\cite{wang2025leveraging} leverages LLMs via Chain-of-Thought for path planning, and T-PRA~\cite{wang2025tunable} employs an LLM-based Actor-Critic framework. However, supervised methods cannot explore paths beyond historical data; heuristic methods greedily optimize local objectives and often yield suboptimal paths; and LLM-based methods incur prohibitive deployment costs.